\documentclass[11pt,a4paper]{article}

\usepackage[a4paper,margin=25mm]{geometry}
\usepackage[T1]{fontenc}
\usepackage{lmodern}
\usepackage{amsmath,amssymb}
\usepackage{bm}
\usepackage{booktabs}
\usepackage[dvipsnames]{xcolor}
\usepackage[authoryear,round]{natbib}
\usepackage[colorlinks=true,linkcolor=blue!50!black,citecolor=blue!50!black,urlcolor=blue!50!black]{hyperref}

\newcommand{\dd}{\mathrm{d}}
\newcommand{\rhat}{\hat{\mathbf{r}}}
\newcommand{\x}{\mathbf{x}}
\newcommand{\uu}{\mathbf{u}}
\newcommand{\grad}{\nabla}
\newcommand{\gradone}{\nabla_{1}}
\newcommand{\Ylm}{Y_{\ell m}}
\newcommand{\code}[1]{\texttt{#1}}
\newcommand{\Ad}{\mathcal{A}_{\ell}}
\newcommand{\sph}{\mathbb{S}^2}

\numberwithin{equation}{section}

\title{\texttt{pyslfp}: theory and numerical methods}
\author{David Al-Attar\\[2pt]
\small Bullard Laboratories, Department of Earth Sciences, University of Cambridge}
\date{October 2026}

\begin{document}
\maketitle

\begin{abstract}
\noindent
This note sets out the physics and numerics behind \texttt{pyslfp}, a Python
library for gravitationally self-consistent sea level fingerprints on a
spherically symmetric, rotating, elastic earth. It states the conventions in
force throughout the code; the generalised load Love numbers, their reciprocity
relations and the axial numbers that enter the rotational feedback; the
spectral-element solution of the radial loading problem at each spherical
harmonic degree; the linear, generalised and nonlinear forms of the sea level
equation together with the iterative methods used to solve them; and the
closed-form treatment of the rotational feedback. A final section describes the
internal-variable formulation of viscoelastic loading that is being added. Every
sign and factor below is the one the code uses.
\end{abstract}

\tableofcontents

%======================================================================
\section{Overview}
\label{sec:overview}
%======================================================================

A change in the distribution of ice, water or sediment on the earth's surface
deforms the solid earth, perturbs its gravitational field and alters its
rotation. Each of these changes the height of the sea surface relative to the
solid surface, and the water so redistributed is itself a load. The resulting
pattern of relative sea level change, the \emph{fingerprint} of the source, is
the solution of the sea level equation
\citep{FarrellClark1976,MitrovicaMilne2003,KendallEtAl2005}. Throughout,
$\xi$ denotes the relative sea level change and $\zeta$ the surface load, a
surface mass density.

\texttt{pyslfp} computes, for a prescribed change in ice thickness (or any
other surface load), the relative sea level change $\xi$, the radial
displacement $u$ of the solid surface, the gravitational potential perturbation
$\phi$ at the surface, and the perturbation $\bm{\omega}$ of the angular velocity.
It does so on a spherically symmetric earth whose response is entirely
characterised, degree by degree, by a table of Love numbers. The library has
three layers.
\begin{enumerate}
\item The radial earth model is held by the separate package
\texttt{planetmodel}, which supplies layered models (density, elastic moduli,
the radii of discontinuities), their gravity, unit systems, and radial meshes
of Gauss--Lobatto--Legendre (GLL) spectral elements.
\item The module \texttt{pyslfp.love\_numbers} solves the linearised,
quasi-static, self-gravitating loading problem of such a model at each degree
$\ell$ (Section~\ref{sec:radial}), and returns a table of generalised Love
numbers (Section~\ref{sec:love}). A default table for PREM without its ocean
layer, for degrees $0\le\ell\le 4096$, is distributed with the library.
\item The sea level solvers in \texttt{pyslfp.core} work entirely in the
spherical harmonic domain on the surface: given the Love numbers they solve
the sea level equation by fixed-point iteration, with the rotational feedback
eliminated exactly within each iteration (Sections~\ref{sec:sle}
and~\ref{sec:rotation}). Two thin drivers sit over this shared engine:
\texttt{pyslfp.elastic}, for which a fingerprint is a single equilibrium
solve, and \texttt{pyslfp.viscoelastic}, whose time stepping drives the same
engine with per-step effective Love numbers (Section~\ref{sec:visco}).
\end{enumerate}
The class \code{EarthModel} binds a Love number table to a set of
non-dimensional scales and spherical harmonic conventions; it adopts the body
(solid surface radius $a$, surface gravity $g$ and gravitational constant $G$)
recorded with the table, so that the sea level equation and its adjoint use
exactly the body for which the Love numbers were computed. A further layer,
\texttt{pyslfp.linear\_operators}, wraps the linear solver as a linear
operator between Hilbert spaces, with its adjoint, for use in inverse problems.

%======================================================================
\section{Conventions and non-dimensionalisation}
\label{sec:conventions}
%======================================================================

\subsection{Spherical harmonics}

Fields on the surface are expanded in real spherical harmonics $\Ylm$ that are
orthonormal on the unit sphere,
\begin{equation}
\int_{\sph} \Ylm\, Y_{\ell' m'}\,\dd\Omega = \delta_{\ell\ell'}\delta_{mm'},
\end{equation}
without the Condon--Shortley phase (in \texttt{pyshtools}, normalisation
\code{'ortho'} and \code{csphase=1}). Orders $m\ge 0$ carry the cosine
harmonics $Y^{c}_{\ell m}\propto P_{\ell m}(\cos\theta)\cos m\varphi$ and $m<0$
the sine harmonics; where the distinction matters we write
$Y^{c}_{\ell m}$ and $Y^{s}_{\ell m}$. The ones needed explicitly are
\begin{equation}
Y_{00} = \frac{1}{\sqrt{4\pi}},\qquad
Y_{20} = \sqrt{\frac{5}{4\pi}}\,\frac{3\cos^{2}\theta-1}{2},\qquad
Y^{c,s}_{21} = \sqrt{\frac{15}{4\pi}}\,\sin\theta\cos\theta\,
\{\cos\varphi,\sin\varphi\}.
\label{eq:Y2}
\end{equation}
A field $f$ on the surface has coefficients $f_{\ell m}$ with
$f=\sum_{\ell m} f_{\ell m}\Ylm$, the series being truncated at a degree $L$
(\code{lmax}). The surface itself is the sphere of radius $b$, the mean radius
of the solid surface; once a Love number table is attached, $b=a$, the radius
of the body for which the table was computed. Surface integrals are therefore
\begin{equation}
\int_{\partial M} f\,\dd S = b^{2}\int_{\sph} f\,\dd\Omega = \sqrt{4\pi}\,b^{2} f_{00},
\label{eq:integral}
\end{equation}
and are evaluated in grid space with the exact quadrature weights of the
Driscoll--Healy or Gauss--Legendre grid in use, which agrees with
\eqref{eq:integral} for band-limited $f$ and is exact also for products such as
$C\xi$ that are not band-limited to $L$.

\subsection{Signs}

The gravitational potential is the physical potential, $\Phi=-G\int\rho/|\x-\x'|\,\dd V'$,
so that $\phi$ is negative near added mass and gravity is $\mathbf{g}=-\grad\Phi$,
with $g=\partial_r\Phi>0$ in the unperturbed state. The centrifugal potential
of rotation with angular velocity $\bm{\Omega}$ is
$-\tfrac12|\bm{\Omega}\times\x|^{2}$ with the same sign convention. Radial
displacement is positive outward.

\subsection{Scales}

The class \code{EarthModelParameters} holds every physical constant in a
non-dimensional form fixed by three base scales: a length $\mathsf{L}$, a
density $\mathsf{P}$ and a time $\mathsf{T}$. These define a mass scale
$\mathsf{M}=\mathsf{P}\mathsf{L}^{3}$, and a quantity with dimensions
$\mathrm{M}^{\alpha}\mathrm{L}^{\beta}\mathrm{T}^{\gamma}$ is stored as its SI
value divided by $\mathsf{M}^{\alpha}\mathsf{L}^{\beta}\mathsf{T}^{\gamma}$.
The same three scales define a \texttt{planetmodel} unit system, so that Love
numbers computed from a model in one set of units can be converted to those of
the sea level solver by the same rule. The default scheme is
\begin{equation}
\mathsf{L} = R_{\oplus},\qquad
\mathsf{P} = \bar\rho = \frac{3M_{\oplus}}{4\pi R_{\oplus}^{3}},\qquad
\mathsf{T} = \sqrt{R_{\oplus}/g_{0}},
\label{eq:scales}
\end{equation}
with $R_{\oplus}=6371$\,km the mean radius, $M_{\oplus}=5.974\times10^{24}$\,kg
and $g_{0}=9.825652323$\,m\,s$^{-2}$. In these units the mean density is one,
the reference gravity is one, and the gravitational constant is
$G\mathsf{P}\mathsf{T}^{2} = 3GM_{\oplus}/(4\pi R_{\oplus}^{2}g_{0})$, close to
$3/4\pi$. The remaining constants are the water and ice densities
$\rho_{w}=1000$ and $\rho_{i}=917$\,kg\,m$^{-3}$, the rotation rate
$\Omega=7.2722\times10^{-5}$\,rad\,s$^{-1}$, and the equatorial and polar moments
of inertia $A=8.0096\times10^{37}$ and $C=8.0359\times10^{37}$\,kg\,m$^{2}$.
The radius $b$, surface gravity $g$ and $G$ are then replaced by those recorded
with the Love number table, where the table records them.

From these the parameters derive five rotational factors, used in
Section~\ref{sec:rotation}:
\begin{gather}
R_{\perp} = \sqrt{\frac{4\pi}{15}}\,\Omega b^{2},\qquad
R_{z} = \sqrt{\frac{16\pi}{45}}\,\Omega b^{2},\qquad
R_{0} = -\sqrt{\frac{16\pi}{9}}\,\Omega b^{2},
\label{eq:rotfactors}\\
I_{\perp} = \sqrt{\frac{5}{12\pi}}\,\frac{\Omega b^{3}}{G(C-A)},\qquad
I_{z} = \sqrt{\frac{5}{9\pi}}\,\frac{\Omega b^{3}}{GC}.
\label{eq:inertiafactors}
\end{gather}
In the code these are \code{rotation\_factor}, \code{axial\_rotation\_factor},
\code{uniform\_rotation\_factor}, \code{inertia\_factor} and
\code{axial\_inertia\_factor}. The $R$ factors map angular velocity to
centrifugal potential coefficients; the $I$ factors map potential coefficients
to angular velocity through MacCullagh's formula.

%======================================================================
\section{Generalised Love numbers}
\label{sec:love}
%======================================================================

\subsection{Load channels}

On a spherically symmetric earth a forcing of degree $\ell$ and order $m$
produces a response of the same degree and order, with amplitudes independent
of $m$. The displacement and potential are written
\begin{equation}
\uu = U(r)\,\Ylm\,\rhat + V(r)\,\gradone\Ylm,\qquad \phi = \phi(r)\,\Ylm,
\label{eq:spheroidal}
\end{equation}
where $\gradone$ is the surface gradient on the unit sphere
\citep{DahlenTromp1998}, so that
\begin{equation}
\int_{\sph}|\gradone\Ylm|^{2}\,\dd\Omega = k^{2}\equiv\ell(\ell+1).
\end{equation}
Following
\citet[eq.~62]{AlAttarEtAl2024}, the action of a surface load on the earth is
split into two parts: a normal traction $-g\zeta\,\rhat$ applied to the
surface, and the gravitational attraction of the surface mass, a source in
Poisson's equation. Each is given its own amplitude, $\zeta^{u}$ for the
traction-like part and $\zeta^{\phi}$ for the potential source, while a true
surface load $\zeta$ acts through both. A third channel, absent from the
paper, is a tangential traction $-g\zeta^{v}\gradone\Ylm$, the tangential
counterpart of the traction part of a load. Finally an external potential
$\psi=(r/a)^{\ell}\Ylm$ represents a harmonic tidal or centrifugal forcing. The
surface response is
\begin{equation}
\begin{aligned}
u_{\ell m} &= h_{\ell}\,\zeta_{\ell m} + h^{u}_{\ell}\,\zeta^{u}_{\ell m}
 + h^{\phi}_{\ell}\,\zeta^{\phi}_{\ell m} + h^{v}_{\ell}\,\zeta^{v}_{\ell m}
 + h^{t}_{\ell}\,\psi_{\ell m},\\
v_{\ell m} &= l_{\ell}\,\zeta_{\ell m} + l^{u}_{\ell}\,\zeta^{u}_{\ell m}
 + l^{\phi}_{\ell}\,\zeta^{\phi}_{\ell m} + l^{v}_{\ell}\,\zeta^{v}_{\ell m}
 + l^{t}_{\ell}\,\psi_{\ell m},\\
\phi_{\ell m} &= k_{\ell}\,\zeta_{\ell m} + k^{u}_{\ell}\,\zeta^{u}_{\ell m}
 + k^{\phi}_{\ell}\,\zeta^{\phi}_{\ell m} + k^{v}_{\ell}\,\zeta^{v}_{\ell m}
 + k^{t}_{\ell}\,\psi_{\ell m},
\end{aligned}
\label{eq:generalised}
\end{equation}
where $u_{\ell m}=U(a)$, $v_{\ell m}=V(a)$ and $\phi_{\ell m}=\phi(a)$ are the
surface values, $\psi_{\ell m}$ is the amplitude of the external potential at
$r=a$, and, because a true load acts through both channels,
\begin{equation}
h_{\ell} = h^{u}_{\ell}+h^{\phi}_{\ell},\qquad
l_{\ell} = l^{u}_{\ell}+l^{\phi}_{\ell},\qquad
k_{\ell} = k^{u}_{\ell}+k^{\phi}_{\ell}.
\end{equation}
The degree $\ell$ is written as a subscript throughout to distinguish it from
the tangential Love number $l$. The twelve numbers
$(h,l,k)^{u,\phi,v,t}_{\ell}$ are what the code computes and stores; the load
sums are derived. They are dimensional: $h^{u,\phi,v}$ and $l^{u,\phi,v}$ are
lengths per unit surface density, $k^{u,\phi,v}$ potentials per unit surface
density, $h^{t}$ and $l^{t}$ lengths per unit potential, and $k^{t}$ is
dimensionless.

The split into channels is not needed to solve the sea level equation forward,
for which $h_{\ell}$ and $k_{\ell}$ suffice. It is needed because the adjoint
of the sea level equation is itself a sea level equation in which an
observation of displacement acts through the traction channel alone and an
observation of potential through the source channel alone
\citep{AlAttarEtAl2024}; similarly an observation of horizontal displacement
requires the tangential channel. Section~\ref{sec:generalised} shows how these
enter the solver.

\subsection{Reciprocity}

The weak form of the loading problem (Section~\ref{sec:radial}) is symmetric.
Each channel enters as a linear functional of the test functions evaluated at
the surface,
\begin{equation}
\zeta^{u}:\ -g a^{2}\,\hat U(a),\qquad
\zeta^{\phi}:\ -a^{2}\,\hat\phi(a),\qquad
\zeta^{v}:\ -g a^{2} k^{2}\,\hat V(a),
\label{eq:loadfunctionals}
\end{equation}
so the matrix of surface responses to the three load channels is symmetric up to
the factors of $g$ and of the norm $k^{2}$ of $\gradone\Ylm$. Writing $\mathsf{A}$
for the discrete system matrix and $\mathbf{e}_{U},\mathbf{e}_{V},\mathbf{e}_{\phi}$
for the surface degrees of freedom, for instance,
$k^{u}_{\ell} = -g a^{2}\,\mathbf{e}_{\phi}^{\mathsf T}\mathsf{A}^{-1}\mathbf{e}_{U}$
and $h^{\phi}_{\ell} = -a^{2}\,\mathbf{e}_{U}^{\mathsf T}\mathsf{A}^{-1}\mathbf{e}_{\phi}$.
This gives the three reciprocity relations
\begin{equation}
g\,h^{\phi}_{\ell} = k^{u}_{\ell},\qquad
h^{v}_{\ell} = k^{2}\,l^{u}_{\ell},\qquad
k^{v}_{\ell} = g\,k^{2}\,l^{\phi}_{\ell},
\label{eq:reciprocity}
\end{equation}
the first being eq.~(64) of \citet{AlAttarEtAl2024}. They hold to rounding in
the computed tables and are checked by \code{LoveNumbers.reciprocity\_residual}.

\subsection{Conventional numbers}

The conventional dimensionless load Love numbers $h'_{\ell},l'_{\ell},k'_{\ell}$
\citep{Farrell1972} are defined against the potential of the surface mass layer
itself. A load $\zeta_{\ell m}\Ylm$ on the sphere of radius $a$ has exterior
potential $-4\pi G a\,\zeta_{\ell m}/(2\ell+1)$ at $r=a$, so
\begin{equation}
h'_{\ell} = \frac{(2\ell+1)\,g}{4\pi G a}\,h_{\ell},\qquad
l'_{\ell} = \frac{(2\ell+1)\,g}{4\pi G a}\,l_{\ell},\qquad
k'_{\ell} = -\frac{2\ell+1}{4\pi G a}\,k_{\ell} - 1,
\label{eq:conventional}
\end{equation}
where $k_{\ell}$ contains the direct attraction of the load and $k'_{\ell}$ does
not. The geodetic tidal numbers are $k^{T}_{\ell}=k^{t}_{\ell}$,
$h^{T}_{\ell}=-g\,h^{t}_{\ell}$ and $l^{T}_{\ell}=-g\,l^{t}_{\ell}$.
In these terms the factor that maps load to sea level (Section~\ref{sec:sle})
is
\begin{equation}
-\Big(h_{\ell}+\frac{k_{\ell}}{g}\Big)
= \frac{4\pi G a}{(2\ell+1)\,g}\,\big(1+k'_{\ell}-h'_{\ell}\big)
= \frac{3}{(2\ell+1)\,\bar\rho}\,\big(1+k'_{\ell}-h'_{\ell}\big),
\label{eq:slcfactor}
\end{equation}
with $\bar\rho = 3g/(4\pi G a)$ the mean density of the body.

The Green's functions for a point mass at angular distance $\theta$ follow from
the addition theorem,
\begin{equation}
G_{u}(\theta) = \sum_{\ell}\frac{2\ell+1}{4\pi a^{2}}\,h_{\ell}\,P_{\ell}(\cos\theta),\qquad
G_{\phi}(\theta) = \sum_{\ell}\frac{2\ell+1}{4\pi a^{2}}\,k_{\ell}\,P_{\ell}(\cos\theta),
\end{equation}
and are summed in the code with the taper $\exp(-10\ell^{2}/L^{2})$.

\subsection{Degrees zero and one}

At $\ell=1$ the problem is posed in the frame of the centre of mass of the
earth and its load. The degree-one exterior potential of a body is
proportional to the offset of its centre of mass from the origin, so in this
frame it vanishes and $\phi(a)=0$ for every forcing. The potential source
channel then does no work, and
\begin{equation}
h^{\phi}_{1}=l^{\phi}_{1}=k^{\phi}_{1}=0,\qquad k_{1}=0,\qquad k'_{1}=-1.
\end{equation}
At $\ell=0$ there is no tangential displacement, and a uniform external
potential is a gauge, so the tidal numbers vanish. The load numbers do not:
conservation of mass fixes the exterior potential of the loaded body,
\begin{equation}
k_{0} = -4\pi G a,\qquad k'_{0}=0,
\end{equation}
while $h_{0}$ is the radial compression of the earth under the weight of a
uniform layer.

\subsection{Axial numbers}
\label{sec:axial}

Five further scalars describe the degree-zero response more completely; they
are needed for the component of the rotational feedback along the rotation
axis (Section~\ref{sec:rotation}). A change $\omega_{z}$ in spin rate has a
centrifugal potential whose spherical mean is proportional to $r^{2}$, not
constant, and so is not harmonic; and any uniform radial deformation changes
the trace of the inertia tensor. Neither effect is contained in the surface
response to a load or to a harmonic tide. We define
\begin{itemize}
\item[] $h_{c}$, $k_{c}$: the surface displacement and potential at degree zero
per unit centrifugal potential $(r/a)^{2}Y_{00}$;
\item[] $m_{u}$, $m_{\phi}$, $m_{c}$: the inertia moments
\begin{equation}
m = \int_{M}\rho\,\uu\cdot\x\,\dd V = \sqrt{4\pi}\int_{0}^{a}\rho\,U\,r^{3}\,\dd r
\label{eq:moment}
\end{equation}
of the degree-zero responses to unit forcing in the traction channel, the
source channel and the centrifugal potential respectively.
\end{itemize}
The moment \eqref{eq:moment} is one quarter of the change in the trace of the
inertia tensor produced by the displacement, since
$\delta\!\int\rho|\x|^{2}\dd V = 2\int\rho\,\uu\cdot\x\,\dd V$ for a
mass-conserving displacement and $\operatorname{tr}\mathbf{I}=2\int\rho|\x|^{2}\dd V$.
The shell theorem gives $k_{c}=0$: a spherically symmetric redistribution of
mass inside the body leaves its exterior potential unchanged.

The centrifugal forcing enters the weak form as the functional
$-\int\rho r^{2}\,\partial_{r}\psi\,\hat U\,\dd r$ with $\psi=(r/a)^{2}$
(Section~\ref{sec:forcings}), that is
$-(2/a^{2})\int\rho r^{3}\hat U\,\dd r = -\big(2/(\sqrt{4\pi}a^{2})\big)\,\hat m$,
where $\hat m$ is the moment \eqref{eq:moment} of the test function. The same
symmetry argument as for \eqref{eq:reciprocity} then gives
\begin{equation}
m_{u} = \tfrac12\sqrt{4\pi}\,g\,a^{4}\,h_{c},\qquad
m_{\phi} = \tfrac12\sqrt{4\pi}\,a^{4}\,k_{c},
\label{eq:axialreciprocity}
\end{equation}
checked by \code{axial\_reciprocity\_residual}; in particular $m_{\phi}=0$, a
surface shell attracting nothing inside it. The self-moment
$m_{c}=-\tfrac12\sqrt{4\pi}a^{2}\,\mathbf{b}_{c}^{\mathsf T}\mathsf{A}^{-1}\mathbf{b}_{c}$
is negative when $\mathsf A$ is positive definite: a positive centrifugal
potential $(r/a)^{2}$ pulls inward. A table that lacks the axial numbers
(as read from an older file) has them undefined, and the sea level solver then
sets the deformation terms of the axial feedback to zero.

\subsection{The table}

\code{LoveNumbers} holds the twelve numbers by degree, the five axial numbers,
and the body $(a,g,G)$, in the units of a given scale system. It is computed by
\code{LoveNumbers.from\_model} from any \texttt{planetmodel} model whose
surface is solid, or read from a plain-text file in SI units whose header
records the columns, the body and the axial numbers. A model frozen at a
frequency $\omega$, whose moduli are complex, yields complex numbers for
frequency-domain viscoelastic work; the sea level solver accepts only real
tables.

%======================================================================
\section{The radial problem at degree \texorpdfstring{$\ell$}{l}}
\label{sec:radial}
%======================================================================

\subsection{The bilinear form}

The Love numbers are surface values of the solution of the linearised,
quasi-static, self-gravitating elastic loading problem on a spherically
symmetric, non-rotating, transversely isotropic earth, reduced by the expansion
\eqref{eq:spheroidal}. The weak form is eq.~(D76) of
\citet{AlAttarTromp2014}, whose isotropic elastic block we write for a
transversely isotropic medium with Love's moduli $A,C,F,L,N$. Define, for the
radial scalars of \eqref{eq:spheroidal},
\begin{equation}
X = r\partial_{r}U,\qquad
Y = 2U - k^{2}V,\qquad
S = U + r\partial_{r}V - V,
\label{eq:XYS}
\end{equation}
and similarly $\hat X,\hat Y,\hat S$ for test functions
$\hat w = (\hat U,\hat V,\hat\phi)$. In solid regions $I_{S}$ the integrand is
\begin{equation}
\begin{aligned}
\mathcal{E}_{S}(w,\hat w) ={}& C X\hat X + (A-N)\,Y\hat Y + F\,(X\hat Y+Y\hat X)
 + k^{2}L\,S\hat S + k^{2}(k^{2}-2)\,N\,V\hat V\\
&+ 4\rho\,(\pi G\rho r - g)\,r\,U\hat U
 + k^{2}\rho g r\,(U\hat V + V\hat U)\\
&+ \rho r^{2}\,(\hat U\,\partial_{r}\phi + U\,\partial_{r}\hat\phi)
 + k^{2}\rho r\,(V\hat\phi + \phi\hat V)
 + \frac{1}{4\pi G}\big(r^{2}\,\partial_{r}\phi\,\partial_{r}\hat\phi + k^{2}\phi\hat\phi\big),
\end{aligned}
\label{eq:solidform}
\end{equation}
with $\rho(r)$ the density and $g(r)$ the gravity of the reference model. Under
$A=C=\kappa+\tfrac43\mu$, $F=\kappa-\tfrac23\mu$ and $L=N=\mu$ the elastic
block becomes
\begin{equation}
\kappa\,(X+Y)(\hat X+\hat Y)
+ \tfrac43\mu\,\big(X-\tfrac12 Y\big)\big(\hat X-\tfrac12\hat Y\big)
+ k^{2}\mu\,S\hat S + k^{2}(k^{2}-2)\,\mu\,V\hat V,
\label{eq:isotropic}
\end{equation}
which is the form of \citet{AlAttarTromp2014}: $X+Y$ is $r$ times the
dilatation, and $X-\tfrac12Y$, $S$ and $V$ are the three deviatoric
strain scalars of the spheroidal field.

In fluid regions $I_{F}$, above degree zero, the quasi-static problem does not
determine the displacement: equilibrium requires only that the perturbed
level surfaces of pressure, density and potential coincide. An equipotential
is displaced radially by $-\phi/g$, so the Eulerian density perturbation is
$\rho_{1} = (\partial_{r}\rho/g)\,\phi$ and the fluid enters through the
potential alone, with integrand
\begin{equation}
\mathcal{E}_{F}(\phi,\hat\phi) = \frac{1}{4\pi G}\big(r^{2}\,\partial_{r}\phi\,\partial_{r}\hat\phi
 + k^{2}\phi\hat\phi\big) + \frac{r^{2}\,\partial_{r}\rho}{g}\,\phi\hat\phi.
\label{eq:fluidform}
\end{equation}
The stratification term is taken as zero where $g=0$. At each fluid--solid
boundary $r=d$ the integration by parts leaves the point term
\begin{equation}
\pm\,\rho_{f}\,d^{2}\big(g\,U\hat U + \phi\,\hat U + U\hat\phi\big)\big|_{r=d},
\label{eq:interface}
\end{equation}
where $\rho_{f}$ is the fluid density at the boundary and the sign is positive
when the fluid lies below (as at the core--mantle boundary) and negative when
it lies above (as at the inner-core boundary), $U$ being the value on the
solid side.

The gravitational field outside the domain is eliminated by Dirichlet-to-Neumann
maps. Outside $r=a$ the potential is $\phi(a)(a/r)^{\ell+1}$, whose energy
$\frac{1}{4\pi G}\int_{a}^{\infty}(r^{2}\phi'^{2}+k^{2}\phi^{2})\dd r$ is
$\frac{(\ell+1)a}{4\pi G}\phi(a)^{2}$. When the domain is truncated below at
$r_{0}>0$ (Section~\ref{sec:submesh}), the potential below is continued as the
regular harmonic $\phi(r_{0})(r/r_{0})^{\ell}$, with energy
$\frac{\ell r_{0}}{4\pi G}\phi(r_{0})^{2}$. The full form is therefore
\begin{equation}
\begin{aligned}
\Ad(w,\hat w) ={}& \int_{I_{S}}\mathcal{E}_{S}\,\dd r + \int_{I_{F}}\mathcal{E}_{F}\,\dd r
 + \sum_{d}(\pm)\rho_{f}d^{2}\big(gU\hat U+\phi\hat U+U\hat\phi\big)\\
&+ \frac{(\ell+1)a}{4\pi G}\phi(a)\hat\phi(a)
 + \frac{\ell r_{0}}{4\pi G}\phi(r_{0})\hat\phi(r_{0}),
\end{aligned}
\label{eq:fullform}
\end{equation}
symmetric in its two arguments. The problem for a forcing $\mathcal{F}$ is to
find $w$ with $\Ad(w,\hat w)=\mathcal{F}(\hat w)$ for all admissible $\hat w$.

\subsection{Forcings}
\label{sec:forcings}

The three load channels are the functionals \eqref{eq:loadfunctionals} per unit
amplitude; a true surface load is the sum of the first two,
$-a^{2}\,[g\hat U(a)+\hat\phi(a)]$, which is eq.~(D79) of
\citet{AlAttarTromp2014}. An external potential $\psi(r)\Ylm$ is moved to the
right-hand side term by term: wherever $\phi$ couples to the displacement or to
the fluid density perturbation in \eqref{eq:solidform}, \eqref{eq:fluidform}
and \eqref{eq:interface} it is replaced by the total potential $\phi+\psi$, the
self-gravitation term $(4\pi G)^{-1}|\grad\phi|^{2}$ being unchanged. This gives
\begin{equation}
\mathcal{F}_{\psi}(\hat w) =
-\int_{I_{S}}\rho\big(r^{2}\,\partial_{r}\psi\,\hat U + k^{2}r\,\psi\,\hat V\big)\dd r
-\int_{I_{F}}\frac{r^{2}\,\partial_{r}\rho}{g}\,\psi\,\hat\phi\,\dd r
-\sum_{d}(\pm)\rho_{f}d^{2}\,\psi(d)\,\hat U(d).
\label{eq:tideforce}
\end{equation}
The harmonic tide is $\psi=(r/a)^{\ell}$. The only non-harmonic case used is
$\psi=(r/a)^{2}$ at $\ell=0$, the spherical mean of the centrifugal potential
of a change in spin rate, which gives $h_{c}$, $k_{c}$ and $m_{c}$. The
constant potential, $\psi=1$, gives zero force.

\subsection{Degrees zero and one}

At $\ell=1$ the surface potential degree of freedom is removed, imposing
$\phi(a)=0$ and with it the centre-of-mass frame; the exterior Dirichlet-to-Neumann
term goes with it. No forcing then acts through the potential source
channel, whose functional is $-a^{2}\hat\phi(a)$.

At $\ell=0$ there is no tangential degree of freedom, and the fluid regions are
treated as elastic media with $\mu=0$ rather than through
\eqref{eq:fluidform}: for a purely radial deformation the fluid's displacement
is determined by its compressibility, and \eqref{eq:isotropic} with $k^{2}=0$
reduces to $\kappa(r\partial_{r}U+2U)^{2}$, the energy of radial compression.
The condition $U(0)=0$ is imposed at the centre.

\subsection{Discretisation}

The radial interval is divided into spectral elements whose edges include every
discontinuity of the model. On each element the unknowns are Lagrange
polynomials on $n$ GLL nodes (\code{ngll}, by default five, so quartic
elements), and all integrals are evaluated by the GLL quadrature on the same
nodes. Radial derivatives are $\partial_{r}=J^{-1}\mathsf{D}$, with
$\mathsf{D}$ the GLL differentiation matrix on the reference element and $J$
the element Jacobian. Solid nodes carry $(U,V,\phi)$, fluid nodes $\phi$ alone,
with $U$ and $V$ at a fluid--solid boundary belonging to the solid side; the
degrees of freedom are numbered node by node, so the global matrix is banded
with half-bandwidth $k_{d}=3(n-1)+2$. Moduli, density and gravity are sampled
at the nodes element by element, so that a node shared by two layers carries
both one-sided values.

The matrix is real symmetric for real moduli and complex symmetric otherwise,
and is stored as an upper band. It is first scaled symmetrically by its
diagonal, $\tilde{\mathsf A}=\mathsf{D}_{A}^{-1/2}\mathsf{A}\mathsf{D}_{A}^{-1/2}$,
which equalises the very different units of the displacement and potential
unknowns. A real system is then factorised by banded Cholesky; the gravitational
terms and the fluid stratification term mean that positive definiteness is not
guaranteed, and the code falls back to banded LU when the factorisation fails.
A complex system is solved by banded LU. For each degree one factorisation
serves all forcings: the three load channels and the tide, and at $\ell=0$ the
centrifugal potential.

\subsection{Sub-mesh and graded mesh}
\label{sec:submesh}

At high degree the solution decays rapidly with depth, roughly as
$(r/a)^{\ell+1}$, and is below a fraction $\varepsilon$ (\code{eps}, by default
$10^{-8}$) of its surface value deeper than
$D_{\ell}=a\,(1-\varepsilon^{1/(\ell+1)})$. Each degree $\ell>0$ is therefore
solved on the sub-mesh of elements above this depth, closed below by the
interior Dirichlet-to-Neumann term of \eqref{eq:fullform}; the displacement and
potential below are set to zero.

A uniform mesh adequate for degree $\ell$ has elements no wider than
$0.1a/(\ell+1)$. At depth $d$ below the surface the highest degree whose
solution reaches that depth satisfies $\ell+1 = \ln\varepsilon/\ln(1-d/a)$, so
an element there need be no narrower than $0.1a\,\ln(1-d/a)/\ln\varepsilon$.
The graded mesh used by default for a table to degree $L$ has element width
\begin{equation}
w(r) = \max\Big(\frac{0.1\,a}{L+1},\ \frac{0.1\,a\,\ln(1-d/a)}{\ln\varepsilon}\Big),
\qquad d=a-r,
\label{eq:graded}
\end{equation}
about $d/184$ at depth for $\varepsilon=10^{-8}$. Elements are laid down from the
top of each layer of the model downwards; the last span above a discontinuity
is divided into one or two equal elements, none wider than $w$ at its top.
Every degree thereby sees at least the resolution the uniform rule would give
it throughout its own sub-mesh, while the whole mesh has a few hundred elements
rather than the $10(L+1)$ of the uniform rule at $L$.

%======================================================================
\section{The sea level equation}
\label{sec:sle}
%======================================================================

\subsection{The linear equation}

Let $C$ be the ocean function of the background state, one over the oceans and
zero elsewhere, $\mathcal{A}_{o}=\int C\,\dd S$ the ocean area, and $\zeta_{d}$ a
prescribed direct load, for example $\rho_{i}\,\Delta I$ for a change $\Delta I$
in grounded ice thickness. With the shorelines held fixed the ocean responds
with a water load $\rho_{w}C\xi$, so the total load is
\begin{equation}
\zeta = \zeta_{d} + \rho_{w}\,C\,\xi.
\label{eq:totalload}
\end{equation}
The sea surface remains an equipotential of the gravitational and centrifugal
potentials, so its height changes by $-(\phi+\psi)/g$ plus a constant, and the
relative sea level change, the change in the height of the sea surface above
the solid surface, is
\begin{equation}
\xi = -u - \frac{\phi+\psi}{g} + c,
\label{eq:sle}
\end{equation}
where $u_{\ell m}=h_{\ell}\zeta_{\ell m}$ and $\phi_{\ell m}=k_{\ell}\zeta_{\ell m}$
apart from the rotational terms of Section~\ref{sec:rotation}, and $c$ is a
spatial constant. The constant is fixed by conservation of mass, the water added
to the oceans being the mass removed from the direct load,
\begin{equation}
\rho_{w}\int C\,\xi\,\dd S = -\int\zeta_{d}\,\dd S
\quad\Longrightarrow\quad
c = \bar\xi_{E} - \frac{1}{\mathcal{A}_{o}}\int C\,\tilde\xi\,\dd S,
\qquad \bar\xi_{E} = -\frac{1}{\rho_{w}\mathcal{A}_{o}}\int\zeta_{d}\,\dd S,
\label{eq:masscons}
\end{equation}
where $\tilde\xi=-u-(\phi+\psi)/g$ and $\bar\xi_{E}$ is the eustatic (barystatic)
sea level change. Equations \eqref{eq:totalload}--\eqref{eq:masscons} are the
linear sea level equation of \citet{FarrellClark1976}, with the rotational terms
of \citet{MilneMitrovica1998}.

\subsection{Fixed-point structure}

Combining \eqref{eq:totalload} and \eqref{eq:sle}, the equation is a fixed point
for the load,
\begin{equation}
\zeta = \mathcal{G}(\zeta) \equiv \zeta_{d} + \rho_{w}\,C\,\mathcal{P}\big[\mathcal{S}\zeta\big],
\qquad
(\mathcal{S}\zeta)_{\ell m} = -\Big(h_{\ell}+\frac{k_{\ell}}{g}\Big)\zeta_{\ell m}
+ (\text{rotation}),
\label{eq:fixedpoint}
\end{equation}
where $\mathcal{P}$ adds the constant \eqref{eq:masscons}. The iteration works
on the coefficients of the load. Each evaluation of $\mathcal{G}$ multiplies
the load coefficients by the combined factor $-(h_{\ell}+k_{\ell}/g)$, replaces
the degree-2 coefficients of orders 0 and 1 by those of
Section~\ref{sec:rotation}, synthesises $\tilde\xi$ on the grid with one inverse
transform, applies the mass-conservation shift using the grid quadrature,
multiplies by $\rho_{w}C$, adds $\zeta_{d}$, and returns to coefficients with
one forward transform. The displacement and potential are synthesised only
once, from the converged load. The starting load is the eustatic one,
$\zeta^{(0)} = \zeta_{d}+\rho_{w}C\,\bar\xi_{E}$.

By \eqref{eq:slcfactor} the linear part of $\mathcal{G}$ is the product of
multiplication by $C$, with norm one on $L^{2}$, and a spherical harmonic
multiplier of size $3\rho_{w}(1+k'_{\ell}-h'_{\ell})/[(2\ell+1)\bar\rho]$, with
the degree-zero part removed by the mass constraint. Since
$\rho_{w}/\bar\rho\approx0.18$ and the multiplier decays as $1/(2\ell+1)$,
the plain Picard iteration contracts, most slowly at the lowest degrees.

\subsection{Anderson acceleration}

The iteration is accelerated by Anderson mixing \citep{WalkerNi2011}. With $x_{k}$ the current
vector of load coefficients, $f_{k}=\mathcal{G}(x_{k})-x_{k}$ its residual, and
the last $p$ differences $\Delta x_{i}=x_{i+1}-x_{i}$, $\Delta f_{i}=f_{i+1}-f_{i}$
gathered as the columns of $\Delta X_{k}$ and $\Delta F_{k}$, the update is
\begin{equation}
\gamma_{k} = \arg\min_{\gamma}\big\|f_{k}-\Delta F_{k}\gamma\big\|_{2},\qquad
x_{k+1} = \mathcal{G}(x_{k}) - \big(\Delta X_{k}+\Delta F_{k}\big)\gamma_{k}.
\label{eq:anderson}
\end{equation}
The least-squares problem is solved through its normal equations
$(\Delta F_{k}^{\mathsf T}\Delta F_{k})\gamma_{k}=\Delta F_{k}^{\mathsf T}f_{k}$,
the $p\times p$ Gram matrix being updated by one row and column per iteration
from ring buffers of the differences, with a minimum-norm least-squares solve
guarding against rank deficiency. The default memory is $p=5$; $p=0$ recovers
the Picard iteration. The iteration stops when the change of the ocean sea level
between iterates satisfies
\begin{equation}
\frac{\max|C\xi^{(k)}-C\xi^{(k-1)}|}{\max|C\xi^{(k)}|}\le\texttt{rtol},
\end{equation}
with $\texttt{rtol}=10^{-9}$ by default.

\subsection{The generalised equation and the adjoint}
\label{sec:generalised}

The generalised sea level equation adds three forcings to the direct load: a
displacement load $\zeta^{u}$ acting through the traction channel, a potential
load $\zeta^{\phi}$ acting through the source channel, and an angular momentum
forcing $\mathbf{L}$ in the rotational equations. The first two contribute static
terms independent of the iteration,
\begin{equation}
u^{s}_{\ell m} = h^{u}_{\ell}\zeta^{u}_{\ell m} + h^{\phi}_{\ell}\zeta^{\phi}_{\ell m},\qquad
\phi^{s}_{\ell m} = k^{u}_{\ell}\zeta^{u}_{\ell m} + k^{\phi}_{\ell}\zeta^{\phi}_{\ell m},
\end{equation}
which are added to $u$ and $\phi$, so that
$\tilde\xi_{\ell m}$ gains $-(u^{s}_{\ell m}+\phi^{s}_{\ell m}/g)$; their
degree-2 and degree-0 coefficients also enter the rotational feedback
(Section~\ref{sec:rotation}). Mass conservation \eqref{eq:masscons} involves
$\zeta_{d}$ alone.

The reason for these channels is the adjoint. The fingerprint operator
$\mathcal{F}:\zeta_{d}\mapsto(\xi,u,\phi,\bm{\omega})$, with the $L^{2}$ inner
product on the sphere of radius $b$ for the fields and the Euclidean one for
$\bm{\omega}$, has an adjoint that is again a solution of the generalised equation
\citep{AlAttarEtAl2024}. Its adjoint applied to
$(\xi^{\dagger},u^{\dagger},\phi^{\dagger},\bm{\omega}^{\dagger})$ is the
relative sea level change of the generalised equation with
\begin{equation}
\zeta_{d} = \xi^{\dagger},\qquad
\zeta^{u} = -u^{\dagger},\qquad
\zeta^{\phi} = -g\,\phi^{\dagger},\qquad
\mathbf{L} = -g\big(\bm{\omega}^{\dagger}-\Psi^{*}\phi^{\dagger}\big),
\label{eq:adjoint}
\end{equation}
where $\Psi:\bm{\omega}\mapsto\psi$ is the centrifugal potential map of
\eqref{eq:centrifugal} below and
\begin{equation}
\Psi^{*}f = b^{2}\big(R_{\perp}f^{c}_{21},\ R_{\perp}f^{s}_{21},\ R_{z}f_{20}+R_{0}f_{00}\big)
\end{equation}
its adjoint. This identity is exact only when the solver uses the reciprocal
numbers of a single table and the body $(a,g,G)$ of that table, which is why
\code{EarthModel} adopts the body from the table.

\subsection{The nonlinear equation with migrating shorelines}

For large changes of ice the ocean function itself changes, and the background
must be carried as fields: the ice thickness $I$ and the sea level $S$, the
height of the sea surface above the solid surface (ocean depth where positive,
minus the topography on land). The ocean function is set by flotation,
\begin{equation}
C = H\big(\rho_{w}S - \rho_{i}I\big),
\label{eq:flotation}
\end{equation}
with $H$ the Heaviside function, so that ice thinner than the water depth that
would float it is part of the ocean, and only grounded ice loads the earth as
ice. Optionally a closed basin (the Caspian) is excluded from $C$. Given an
initial state $(I_{0},S_{0},C_{0})$, an ice thickness change $\Delta I$, a
sediment thickness change $\Delta H_{s}$ of density $\rho_{s}$ (2300\,kg\,m$^{-3}$
by default) and a dynamic sea level change $\xi_{\mathrm{dyn}}$, the final ice
is $I=I_{0}+\Delta I$ and the solver iterates on $(S_{n},C_{n})$ from
$(S_{0},C_{0})$. Each iteration forms the load
\begin{equation}
\zeta_{n} = \rho_{i}\big[(1-C_{n})I - (1-C_{0})I_{0}\big]
 + \rho_{w}\big(C_{n}S_{n} - C_{0}S_{0}\big) + \rho_{s}\Delta H_{s},
\end{equation}
computes $\tilde\xi_{n}=\mathcal{S}\zeta_{n}$ as in \eqref{eq:fixedpoint},
including the rotational terms but without the constant, and the provisional sea
level $\hat S = S_{0}+\tilde\xi_{n}-\Delta H_{s}+\xi_{\mathrm{dyn}}$, sediment
raising the solid surface. The constant is a eustatic shift that conserves the
total of water, grounded ice and sediment over the current ocean,
\begin{gather}
\Delta_{n} = \frac{M_{n} - \rho_{w}\int C_{n}\hat S\,\dd S}{\rho_{w}\int C_{n}\,\dd S},\\
M_{n} = \rho_{w}\!\int C_{0}S_{0}\,\dd S
 - \int\rho_{i}\big[(1-C_{n})I-(1-C_{0})I_{0}\big]\dd S
 - \int\rho_{s}\Delta H_{s}\,\dd S,
\end{gather}
after which
\begin{equation}
\xi_{n+1} = \tilde\xi_{n}+\Delta_{n},\qquad
S_{n+1} = \hat S + \Delta_{n},\qquad
C_{n+1} = H\big(\rho_{w}S_{n+1}-\rho_{i}I\big).
\end{equation}
This is the formulation of \citet{KendallEtAl2005}. The iteration is plain
Picard, stopped when $\max|\xi_{n+1}-\xi_{n}|\le\texttt{rtol}\,\max|\xi_{n}|$;
the error is measured on the whole grid, not only over the ocean. The result is a new state
$(I,S_{n+1})$ together with $\xi$, $u$, $\phi$ and $\bm{\omega}$, the last three
synthesised from the load of the final iteration.

%======================================================================
\section{Rotational feedback}
\label{sec:rotation}
%======================================================================

\subsection{Centrifugal potential}

Let the angular velocity be $\Omega\hat{\mathbf z}+\bm{\omega}$. To first order
the centrifugal potential changes by
\begin{equation}
\psi = -\Omega\omega_{z}r^{2} + \Omega z\,(\bm{\omega}\cdot\x)
= \Omega r^{2}\Big[\sin\theta\cos\theta\,(\omega_{x}\cos\varphi+\omega_{y}\sin\varphi)
+ \omega_{z}\big(\cos^{2}\theta-\tfrac13\big) - \tfrac23\,\omega_{z}\Big].
\end{equation}
With \eqref{eq:Y2}, its coefficients on the surface $r=b$ are
\begin{equation}
\psi^{c}_{21} = R_{\perp}\,\omega_{x},\qquad
\psi^{s}_{21} = R_{\perp}\,\omega_{y},\qquad
\psi_{20} = R_{z}\,\omega_{z},\qquad
\psi_{00} = R_{0}\,\omega_{z},
\label{eq:centrifugal}
\end{equation}
with the factors \eqref{eq:rotfactors}. The degree-2 parts are of the form
$(r/b)^{2}Y_{2m}$, harmonic, and act on the earth through the tidal Love numbers
$h^{t}_{2}$, $k^{t}_{2}$. The degree-0 part is $\psi_{00}(r/b)^{2}Y_{00}$, not
harmonic, and acts through the axial numbers $h_{c}$, $k_{c}$ and $m_{c}$.
The \emph{transverse} components $\omega_{x},\omega_{y}$ (true polar wander)
therefore couple to degree 2, order 1, and the \emph{axial} component
$\omega_{z}$ (change in length of day) to degree 2, order 0, and degree 0.

\subsection{Liouville equation and MacCullagh's formula}

On the time scales of loading the Chandler wobble is not excited and the
linearised Liouville equation reduces to its quasi-static balance. With
$\Delta\mathbf{I}$ the perturbation of the inertia tensor and $\mathbf{h}$ a
relative angular momentum, it reads
\begin{equation}
(C-A)\,\omega_{x} = \Omega\,\Delta I_{xz} + h_{x},\qquad
(C-A)\,\omega_{y} = \Omega\,\Delta I_{yz} + h_{y},\qquad
C\,\omega_{z} = -\Omega\,\Delta I_{zz} - h_{z}.
\label{eq:liouville}
\end{equation}
The inertia perturbation is read from the degree-2 potential by MacCullagh's
formula. Expanding $-G\int\rho/|\x-\x'|$ to second order outside the mass and
writing $\Delta\mathbf{I}$ with $\Delta T=\operatorname{tr}\Delta\mathbf{I}$,
\begin{equation}
\phi^{c}_{21} = \frac{3G}{b^{3}}\sqrt{\frac{4\pi}{15}}\,\Delta I_{xz},\qquad
\phi^{s}_{21} = \frac{3G}{b^{3}}\sqrt{\frac{4\pi}{15}}\,\Delta I_{yz},\qquad
\phi_{20} = \frac{G}{b^{3}}\sqrt{\frac{\pi}{5}}\,\big(3\Delta I_{zz}-\Delta T\big).
\label{eq:maccullagh}
\end{equation}
Eliminating $\Delta\mathbf{I}$ between \eqref{eq:liouville} and
\eqref{eq:maccullagh} gives, with the factors \eqref{eq:inertiafactors},
\begin{equation}
\omega_{x,y} = I_{\perp}\,\phi^{c,s}_{21} + \frac{h_{x,y}}{C-A},\qquad
\omega_{z} = -I_{z}\,\phi_{20} - \frac{\Omega}{3C}\,\Delta T - \frac{h_{z}}{C}.
\label{eq:omega}
\end{equation}
The transverse equations involve only the degree-2 potential. The axial
equation also involves the trace $\Delta T$, which the exterior potential does
not see. The angular momentum forcing of the generalised equation enters as
$\mathbf{h}=-\mathbf{L}$; that is, $\mathbf{L}$ adds
$-\mathbf{L}_{x,y}/(C-A)$ to $\omega_{x,y}$ and $+L_{z}/C$ to $\omega_{z}$.

\subsection{The trace}

The change in trace has two sources. A surface mass layer $\sigma$ on the
sphere of radius $b$ contributes $2b^{2}\int\sigma\,\dd S = 2\sqrt{4\pi}\,b^{4}\sigma_{00}$,
and the mass-like sources are the true load and the potential load, so
$\sigma_{00}=\zeta_{00}+\zeta^{\phi}_{00}$; the displacement load is a traction
and carries no mass. The deformation contributes $4m$, with $m$ the moment
\eqref{eq:moment} of the degree-0 displacement, which by
Section~\ref{sec:axial} is
\begin{equation}
m = (m_{u}+m_{\phi})\,\zeta_{00} + m_{u}\,\zeta^{u}_{00} + m_{\phi}\,\zeta^{\phi}_{00} + m_{c}\,\psi_{00}.
\end{equation}
Hence
\begin{equation}
\frac{\Omega}{3C}\,\Delta T
= \frac{4\Omega}{3C}\,m + \frac{2\Omega\sqrt{4\pi}\,b^{4}}{3C}\,\big(\zeta_{00}+\zeta^{\phi}_{00}\big).
\end{equation}

\subsection{Closed-form fixed points}

The potential in \eqref{eq:omega} contains the response to the centrifugal
potential, $\phi_{2m}=\phi^{L}_{2m}+k^{t}_{2}\psi_{2m}$, where
$\phi^{L}_{2m} = k_{2}\zeta_{2m}+\phi^{s}_{2m}$ is the potential of the loads.
Likewise the moment depends on $\psi_{00}=R_{0}\omega_{z}$. Since
\eqref{eq:centrifugal} is linear in $\bm{\omega}$, each component is a scalar
fixed point, solved in closed form:
\begin{gather}
\omega_{x,y} = \frac{I_{\perp}\,\phi^{L,c/s}_{21} - L_{x,y}/(C-A)}
{1 - I_{\perp}\,k^{t}_{2}\,R_{\perp}},
\label{eq:transverse}\\[4pt]
\omega_{z} = \frac{-I_{z}\,\phi^{L}_{20}
 - \dfrac{4\Omega}{3C}\big[(m_{u}+m_{\phi})\zeta_{00}+m_{u}\zeta^{u}_{00}+m_{\phi}\zeta^{\phi}_{00}\big]
 - \dfrac{2\Omega\sqrt{4\pi}\,b^{4}}{3C}\big(\zeta_{00}+\zeta^{\phi}_{00}\big)
 + L_{z}/C}
{1 + I_{z}\,k^{t}_{2}\,R_{z} + \dfrac{4\Omega}{3C}\,m_{c}\,R_{0}}.
\label{eq:axialomega}
\end{gather}
In the axial denominator both corrections are positive ($k^{t}_{2}>0$, and
$m_{c}<0$, $R_{0}<0$), so the feedback of the axial deformation on the spin
rate is stabilising. In the linear equation $\zeta_{00}$ vanishes at convergence
by mass conservation, so the degree-0 terms in $\zeta$ matter chiefly for the
generalised equation and for loads that do not conserve mass.

Given $\bm{\omega}$, the rotational contributions to the fields are
\begin{equation}
u_{2m} \mathrel{+}= h^{t}_{2}\psi_{2m},\quad
\phi_{2m} \mathrel{+}= k^{t}_{2}\psi_{2m},\quad
u_{00} \mathrel{+}= h_{c}\psi_{00},\quad
\phi_{00} \mathrel{+}= k_{c}\psi_{00},
\end{equation}
for the degree-2 coefficients of orders 0 and 1 and the degree-0 coefficient,
and the relative sea level coefficients become
\begin{equation}
\tilde\xi_{2m} = -\Big(u_{2m}+\frac{\phi_{2m}+\psi_{2m}}{g}\Big),\qquad
\tilde\xi_{00} \mathrel{-}= \Big(h_{c}+\frac{k_{c}+1}{g}\Big)\psi_{00}.
\end{equation}
The last is uniform and is absorbed by the mass-conservation constant; it is
kept so that $\tilde\xi$ is the complete local response. The returned potential
$\phi$ is the gravitational perturbation without $\psi$.

Because the rotational fixed point is solved exactly inside every evaluation of
$\mathcal{G}$, the angular velocity, displacement and potential returned are all
consistent with the same load, and the outer iteration has only the ocean
loading to converge.

\subsection{Rotational stability and the choice of moments}

The equilibrium moments of inertia in \eqref{eq:transverse} and
\eqref{eq:axialomega} are specified quantities: they cannot be computed from
the spherically symmetric model in which the elastic problem is solved, whose
inertia tensor is isotropic \citep{YuEtAl2025}. The transverse denominator
$1-I_{\perp}k^{t}_{2}R_{\perp}$ equals $1-k^{t}_{2}/k_{\mathrm{sus}}$, with
$k_{\mathrm{sus}}=3G(C-A)/(\Omega^{2}b^{5})$ the oblateness the specified
moments sustain, and it must remain positive for every tidal response the
model can produce. The largest such response is the hydrostatic fluid limit
$k_{\mathrm{hyd}}$, the degree-2 tidal number of the model with its rigidity
set to zero, which depends only on the density profile (Kelvin's
$3/(2(\ell-1))$ for a homogeneous body); the library computes it directly as
the zero-rigidity solid rather than as the $s\to0$ limit of a viscoelastic
$k^{t}_{2}(s)$, which degrades numerically \citep{MitrovicaEtAl2005}. For the
Earth the observed moments satisfy $k_{\mathrm{hyd}}<k_{\mathrm{sus}}$, the
excess being the non-hydrostatic part of the figure, and the specified
observed values are the correct default. For any other body the moments may
instead be constructed from the model itself in the manner of
\citet{MitrovicaEtAl2005}: $C-A=(1+\delta)\,k_{\mathrm{hyd}}\Omega^{2}b^{5}/(3G)$
with $\delta$ a prescribed non-hydrostatic fraction ($0.8$ per cent for the
Earth), which makes the margin $\delta/(1+\delta)$ whatever the rheology.
A configuration with $k_{\mathrm{hyd}}\ge k_{\mathrm{sus}}$ has a transverse
feedback on the wrong side of the secular resonance and is refused by the
time-domain solvers; the axial feedback is unconditionally stable, both
denominator corrections in \eqref{eq:axialomega} being positive, and may be
retained alone.

%======================================================================
\section{Viscoelastic extension (in development)}
\label{sec:visco}
%======================================================================

This section describes a formulation being implemented in
\texttt{pyslfp.gia}; its interface is not yet stable.

\subsection{Rheology and internal variables}

The extension treats a generalised Maxwell body with branches
$k=1,\dots,K$ of shear modulus $\mu_{k}$ and relaxation time $\tau_{k}$, and a
long-term shear modulus $\mu_{\infty}$, in the internal-variable formulation of
\citet{AlAttarTromp2014} (see also \citealp{CrawfordEtAl2017}). The same
formulation underlies the time-domain treatment of glacial isostatic adjustment,
and of its sensitivity kernels by the adjoint method, in laterally
heterogeneous earth models by \citet{CrawfordEtAl2018}; here the earth is
spherically symmetric and each degree decouples. With $d$ the
deviatoric strain, each branch carries a trace-free symmetric internal variable
$m_{k}$, and
\begin{equation}
\mathbf{T} = \kappa\,(\operatorname{tr}\bm{\varepsilon})\,\mathbf{1}
 + 2\mu_{\infty}d + \sum_{k}2\mu_{k}\,(d-m_{k}),\qquad
\dot m_{k} = \frac{d-m_{k}}{\tau_{k}}.
\label{eq:maxwell}
\end{equation}
The unrelaxed shear modulus is $\mu_{U}=\mu_{\infty}+\sum_{k}\mu_{k}$, which is
the elastic modulus of the material. The classical Maxwell body is one branch
with $\mu_{\infty}=0$ and $\tau=\eta/\mu$. The viscosity may be prescribed
directly, or through the relaxation time $\tau(r)$, in which case
$\eta=\mu\tau$ inherits its depth dependence from the shear modulus. Bulk
relaxation is not modelled, fluid regions carry no branches, and branches
require an isotropic material.

\subsection{The degree-basis scalars}

For spheroidal forcing the internal variable of each branch reduces, at each
$(\ell,m)$, to three radial scalars \citep[eqs D68--D73]{AlAttarTromp2014}: $M$,
the purely tangential component, $R$, the radial--radial component, and $S$,
the radial--tangential component; the toroidal components remain zero. They are
driven by the deviatoric strain scalars of \eqref{eq:isotropic},
\begin{equation}
e_{M} = V,\qquad
e_{R} = r\partial_{r}U - U + \tfrac12 k^{2}V = X-\tfrac12 Y,\qquad
e_{S} = r\partial_{r}V - V + U = S,
\end{equation}
branch by branch,
\begin{equation}
\dot M_{k} = \frac{e_{M}-M_{k}}{\tau_{k}},\qquad
\dot R_{k} = \frac{e_{R}-R_{k}}{\tau_{k}},\qquad
\dot S_{k} = \frac{e_{S}-S_{k}}{\tau_{k}},
\end{equation}
the strain scalars being evaluated at the GLL nodes by collocation. The weights
with which they enter the $\mu$-part of \eqref{eq:isotropic} are $\tfrac43$,
$k^{2}$ and $k^{2}(k^{2}-2)$ for $R$, $S$ and $M$.

\subsection{Quasi-static equilibrium and the memory force}

At fixed internal variables the displacement and potential satisfy the elastic
problem with the unrelaxed moduli, forced by the loads and by a memory force:
\begin{gather}
\Ad^{U}(w,\hat w) = \mathcal{F}(\hat w) + b_{\mathrm{mem}}(\hat w),\\
b_{\mathrm{mem}}(\hat w) = \sum_{k}\int\mu_{k}\Big[\tfrac43 R_{k}\,e_{R}(\hat w)
 + k^{2}S_{k}\,e_{S}(\hat w) + k^{2}(k^{2}-2)\,M_{k}\hat V\Big]\dd r,
\end{gather}
where $\Ad^{U}$ is \eqref{eq:fullform} assembled with $\mu_{U}$. This is
the one-sided form of the $\mu$-part of \eqref{eq:isotropic} applied to the
internal variables, and it follows from \eqref{eq:maxwell} by moving the
$-2\mu_{k}m_{k}$ stress to the right-hand side. At $\ell=0$ only the $R$ term
survives; at $\ell=1$ the $M$ term vanishes although $M$ still evolves. The
state of the system is the internal variables alone: the displacement,
potential and, once coupled, the sea level are slaved to them through degree
solves, each a banded back-substitution with the orders as right-hand-side
columns.

\subsection{Time stepping}

With $h_{k}=\Delta t/\tau_{k}$, each scheme eliminates the new internal variable
in favour of the new strain, which leaves an elastic solve with an effective
shear modulus and a memory force built from known quantities.
\begin{itemize}
\item[] \emph{Backward Euler}: $m^{n+1}=(m^{n}+h\,d^{n+1})/(1+h)$, so that
\begin{equation}
\mu_{\mathrm{eff}} = \mu_{\infty}+\sum_{k}\frac{\mu_{k}}{1+h_{k}},\qquad
\tilde m_{k} = \frac{m_{k}^{n}}{1+h_{k}}.
\end{equation}
\item[] \emph{Exponential trapezoid}: with $e=\exp(-h)$, $\varphi_{1}=(1-e)/h$,
$\alpha=\varphi_{1}-e$ and $\beta=1-\varphi_{1}$, the exact update for a strain
linear in time over the step is $m^{n+1}=e\,m^{n}+\alpha\,d^{n}+\beta\,d^{n+1}$,
so that
\begin{equation}
\mu_{\mathrm{eff}} = \mu_{\infty}+\sum_{k}\varphi_{1,k}\,\mu_{k},\qquad
\tilde m_{k} = e_{k}\,m_{k}^{n}+\alpha_{k}\,d^{n};
\end{equation}
$\varphi_{1}$ is evaluated by its Taylor series for small $h$.
\item[] \emph{Exponential Euler}: $m^{n+1}=e\,m^{n}+(1-e)\,d^{n}$, explicit, the
solve at $t^{n}$ using the unrelaxed moduli and $b_{\mathrm{mem}}(m^{n})$.
\end{itemize}
In each case the step's equilibrium is $\Ad^{\mathrm{eff}}(w,\hat w)=\mathcal{F}(\hat w)
+b_{\mathrm{mem}}(\tilde m)$, after which $m^{n+1}$ follows from the update
formula. An elastic region is the limit $\tau\to\infty$: $h=0$, $e=1$,
$\alpha=\beta=0$, $\varphi_{1}=1$, and $m$ remains zero. Backward Euler is first
order and the exponential trapezoid second order; the exponential schemes are
exact for relaxation under strain prescribed piecewise constant (exponential
Euler) or piecewise linear (trapezoid) in time, for any $\Delta t$. One
factorisation per degree and per $(\text{scheme},\Delta t)$ serves every step.

\subsection{Coupling to the sea level equation}

With the effective moduli the instantaneous response of a step to a surface
load is described by per-step effective Love numbers, computed by the elastic
machinery of Section~\ref{sec:radial} on the effective material, together with a
memory field, the surface response $(U,\phi)$ to $b_{\mathrm{mem}}(\tilde m)$
under zero load. The step's sea level equation is then the generalised linear
equation of Section~\ref{sec:sle} with the effective numbers in place of
$h_{\ell},k_{\ell}$ and the memory field as a static term, closed by the same
iteration for the ocean and rotational feedbacks. The converged load gives
$d^{n+1}$ at the nodes and completes the update of the internal variables. In
the limit $\tau_{k}\to\infty$ for every branch the formulation reduces to the
elastic one. The same body can also be treated in the frequency domain, through
the complex moduli of a model frozen at a frequency
(Section~\ref{sec:love}); by the correspondence principle
\citep{WuPeltier1982} the two descriptions are related by a Laplace
transform, which provides an independent check of the time stepping.

\bibliographystyle{plainnat}
\bibliography{pyslfp_theory}

\end{document}
